\chapter{Fazit und Ausblick}

Dieses Kapitel bündelt die wesentlichen Erkenntnisse der Arbeit, bewertet die Erreichung des Forschungsziels und leitet daraus priorisierte Schritte für Forschung und Entwicklung ab. Aussagen über den final integrierten Prototyp stützen sich auf die vollständige End-to-End-Evaluation vom 23.~August~2026. Ergebnisse älterer Systemstände werden ausschließlich als historische Vergleichs- oder Komponentenbefunde herangezogen. Diese Trennung ist notwendig, weil sich unter anderem Korpus, Antwortmodell, CRM-Integration, Reranker und Sicherheitskonfiguration im Verlauf der Entwicklung verändert haben.

\section{Zusammenfassung der Arbeit}

Ausgangspunkt der Arbeit war die Frage, wie ein KI-basiertes Assistenzsystem versicherungsbezogene Anfragen beantworten kann, wenn dafür sowohl strukturierte Kunden-, Policen- und Schadendaten als auch unstrukturierte Versicherungsdokumente benötigt werden. Die besondere Schwierigkeit liegt nicht allein im Auffinden relevanter Informationen. Ebenso wichtig sind die korrekte Zuordnung zu Personen, Verträgen und Produkten, die erkennbare Herkunft der verwendeten Evidenz, der Schutz personenbezogener Daten und ein kontrolliertes Verhalten bei unzureichenden oder unzulässigen Anfragen. Vor diesem Hintergrund wurde eine agentische CRM--RAG-Architektur als wissenschaftlicher Proof of Concept entworfen, implementiert und mehrdimensional evaluiert.

Der entwickelte Prototyp verbindet EspoCRM mit einem dokumentenbasierten RAG-System. Ein deterministisches Top-Level-Routing ordnet Anfragen den Pfaden CRM-only, RAG-only, Combined oder Denied zu. Der CRM-only-Pfad verwendet ausschließlich begrenzte read-only-Werkzeuge und formatiert die ermittelten Fakten ohne abschließende LLM-Generierung. Der RAG-only-Pfad kombiniert BM25, semantische Vektorsuche, Reciprocal Rank Fusion und Cross-Encoder-Reranking, bevor eine quellengebundene Antwort erzeugt wird. Im Combined-Pfad werden strukturierte CRM-Evidenz und Dokumentpassagen getrennt ermittelt, mit Herkunftsmetadaten versehen und erst danach gemeinsam für die Antwortgenerierung bereitgestellt. Eingangsprüfungen, PII-Behandlung, Guardrails, Groundedness-Prüfung, Fallbacks sowie Audit- und Diagnosedaten begrenzen und dokumentieren die Verarbeitung.

Methodisch folgt die Arbeit einem Design-Science-Ansatz. Die zentrale Forschungsfrage wurde durch die vier Qualitätsdimensionen fachliche Korrektheit und Vollständigkeit, quellenbasierte Nachvollziehbarkeit, Sicherheit sowie praktische Einsetzbarkeit operationalisiert. Die Evaluation verbindet Komponenten-, Integrations- und End-to-End-Tests. Dabei wurden Routing und operativer Abschluss bewusst von Retrieval-Erfolg, Antwortqualität, Quellenabdeckung und Groundedness getrennt. Diese Trennung verhindert, dass ein technisch erfolgreich abgeschlossener Fall fälschlich als fachlich korrekte Antwort interpretiert wird.

Im aktuellen vollständigen Referenzlauf wurden 80 eindeutige Fälle mit jeweils 20 Anfragen pro Top-Level-Pfad neu über die reale API verarbeitet. Routing Accuracy, Macro-F1, Strict End-to-End Success und Downstream Operational Completion erreichten jeweils 100~\%; Systemfehler, Timeouts und HTTP-5xx-Antworten traten nicht auf. Das davon getrennte automatisierte, referenzbasierte Qualitäts-Gate bestanden 73 von 80 Fällen beziehungsweise 91{,}25~\%. Der mittlere CRM Fact Recall betrug 97{,}50~\%, der mittlere Requirement Recall 85{,}00~\%. In 39 von 40 dokumentbasierten Fällen befand sich die erwartete Datei sowohl im Reranker-Ergebnis als auch im Generierungskontext. Die Expected-Source Citation Rate lag jedoch nur bei 87{,}50~\%, und die mittlere Claim Support Rate erreichte 65{,}32~\%. Von 60 nicht abgelehnten Fällen überschritten 56 die Groundedness-Schwelle.

Die Fehleranalyse zeigt damit einen klaren Übergabeverlust zwischen gefundener Evidenz und finaler Antwort. Nur einer der sieben nicht bestandenen Fälle war primär auf das Verfehlen der erwarteten Datei im Retrieval zurückzuführen. In mehreren anderen Fällen war die relevante Evidenz vorhanden, wurde aber nicht vollständig in Muss-Anforderungen, Aussagen und Zitationen überführt oder ging im Fallback verloren. Die stärkste Evidenz der Arbeit betrifft daher die technische Integration, die korrekte Pfadauswahl und den begrenzten CRM-Zugriff. Für fachliche Vollständigkeit, Claim-Unterstützung und Zitationsabdeckung bleibt trotz der deutlichen Verbesserung weiterer Entwicklungs- und Validierungsbedarf.

\section{Wissenschaftliche und technische Beiträge}

Der wissenschaftlich-konzeptionelle Beitrag liegt in einer kontrollierten Form agentischer Evidenzorchestrierung. Die Architektur entscheidet nicht frei über beliebige Werkzeuge, sondern wählt deterministisch zwischen vier klar abgegrenzten Verarbeitungspfaden. Agentik bezeichnet damit die koordinierte, zustandsabhängige Nutzung begrenzter Komponenten und nicht die autonome Planung eines frei agierenden Sprachmodells. Diese Ausgestaltung reduziert Flexibilität zugunsten von Nachvollziehbarkeit, Reproduzierbarkeit und Zugriffskontrolle und ist insbesondere für sensible Unternehmensdaten geeignet.

Ein zweiter Beitrag ist die herkunftserhaltende Verbindung strukturierter und unstrukturierter Evidenz. Individuelle CRM-Fakten und allgemeine Dokumentaussagen bleiben bis zur Ausgabe unterscheidbar. Dadurch kann das System beispielsweise einen im CRM gespeicherten Selbstbehalt nennen, ohne aus diesem Feld allein eine allgemeine Deckungszusage abzuleiten. Umgekehrt darf eine allgemeine Versicherungsbedingung nicht als kundenspezifisches Vertragsmerkmal erscheinen. Die deterministische Validierung von Kontakt-, Policen- und Schadenbeziehungen sowie die Auswahl aktueller Policen ergänzen semantische Verfahren dort, wo zeitliche und relationale Geschäftslogik erforderlich ist.

Der methodische Beitrag besteht in einem Evaluationsrahmen, der technische Ausführbarkeit, fachliche Anforderungen und Sicherheitsentscheidungen als getrennte Gates behandelt. Routing Accuracy, Strict End-to-End Success, Downstream Operational Completion und Automated Overall Pass messen unterschiedliche Sachverhalte und werden nicht zu einer künstlichen Gesamtkennzahl verdichtet. Ergänzende Metriken für CRM-Fakten, erwartete Dokumente, Requirements, Claims, Zitationen und Groundedness machen sichtbar, an welcher Stelle der Verarbeitungskette Qualität verloren geht. Der versionierte Vergleich unterschiedlicher Systemstände zeigt außerdem, weshalb Ergebnisse nur zusammen mit Datensatz, Modell-, Prompt- und Konfigurationsstand interpretiert werden dürfen.

Der technische Beitrag umfasst eine lauffähige Integration von FastAPI, EspoCRM, persistentem MCP-Zugriff, Chroma, hybridem Retrieval, Cross-Encoder-Reranking, externer Antwortgenerierung, NeMo Guardrails und einer faktensensitiven Groundedness-Prüfung. Die fünf CRM-Werkzeuge sind ausschließlich lesend; Schreib-, Lösch-, Export- und unbeschränkte Listenoperationen stehen dem System nicht zur Verfügung. Readiness-Prüfungen, Auditdaten und instrumentierte Stage-Latenzen machen nicht nur das Ergebnis, sondern auch die tatsächlich ausgeführten Komponenten nachvollziehbar.

Der empirische Beitrag liegt in der gemeinsamen Ausführung und automatisierten Referenzbewertung aller vier Verarbeitungspfade im finalen Systemstand. Darüber hinaus zeigen die Komponentenbefunde, dass MiniLM im untersuchten CPU-Setup mit Top-1 = 0{,}969, MRR@5 = 0{,}982 und nDCG@5 = 0{,}986 bei einer Medianlatenz von 467{,}8~ms einen günstigen Qualitäts-Latenz-Kompromiss bietet. Der Beitrag der Arbeit besteht somit nicht in einem neuen Sprach-, Retrieval- oder Reranking-Algorithmus. Er liegt in der domänenspezifischen Integration, der expliziten Evidenzherkunft, der Begrenzung von Werkzeugen und der gemeinsamen, mehrdimensionalen Evaluation dieser Elemente.

\section{Gesamtbewertung der Zielerreichung}

Die zentrale Forschungsfrage kann auf Ebene des untersuchten Proof of Concept positiv, jedoch nur unter klaren Bedingungen beantwortet werden. Eine agentische RAG-Architektur für den Versicherungsbereich kann fachlich kontrollierbarer, quellenbezogen nachvollziehbarer, sicherer und praktisch nutzbarer gestaltet werden, wenn Datenquellen anfrageabhängig ausgewählt, Werkzeugrechte begrenzt, Evidenzarten getrennt gekennzeichnet und technische, fachliche sowie sicherheitsbezogene Freigaben unabhängig voneinander geprüft werden. Die Ergebnisse belegen, dass diese Prinzipien gemeinsam in einer lauffähigen Architektur umgesetzt werden können.

Hinsichtlich der fachlichen Korrektheit und Vollständigkeit ist das Ziel teilweise erreicht. Die hohen Werte für CRM Fact Recall und erwartete Dokumentpräsenz zeigen, dass die benötigten Informationen meist verfügbar waren. Der Requirement Recall von 85{,}00~\% und die sieben nicht bestandenen Fälle belegen jedoch, dass verfügbare Evidenz nicht immer vollständig in die Antwort übernommen wurde. Groundedness allein schließt diese Lücke nicht, weil sie die Stützung vorhandener Aussagen und nicht das Fehlen erforderlicher Aussagen bewertet.

Die quellenbasierte Nachvollziehbarkeit ist technisch umgesetzt, aber noch nicht lückenlos. Dokumentquellen werden mit Datei- und Seitenbezug ausgegeben, CRM-Fakten bleiben als eigene Evidenzart erkennbar, und die Präzision der zurückgegebenen Quellenlinks war im aktuellen Lauf hoch. Gleichzeitig lagen erwartete Quellen nicht in allen dokumentbasierten Antworten vor, und der Claim-Coverage-Proxy sowie die Claim Support Rate zeigen verbleibende Beleglücken. Sichtbare Quellen dürfen deshalb nicht mit vollständiger Aussagenstützung gleichgesetzt werden.

Für Sicherheit und praktische Einsetzbarkeit liefert der Prototyp einen belastbaren technischen, aber keinen produktiven Nachweis. Alle 20 Denied-Fälle des aktuellen Laufs wurden sicher abgeschlossen, ohne dass fachliche Datenquellen aufgerufen wurden. Die read-only-Ausgestaltung reduziert mögliche Schäden auch dann, wenn eine sprachliche Anfrage fehlerhaft klassifiziert würde. Gleichzeitig fehlen für reale Kundendaten weiterhin eine eigene API-Authentifizierung, rollenbasierte Autorisierung, aktiviertes TLS, eine risikoadäquate Fail-closed-Strategie und produktionsgerechtes Monitoring. Der sequenzielle 80-Fall-Lauf erlaubt zudem keine Aussage über Parallelität, Skalierbarkeit, Kosten oder Langzeitstabilität.

Das Hauptziel der Arbeit ist damit auf Proof-of-Concept-Ebene erreicht: Die entworfene CRM--RAG-Architektur ist vollständig implementiert, in allen vorgesehenen Pfaden ausführbar und anhand mehrerer Qualitätsdimensionen evaluiert. Nicht erreicht sind eine unabhängige fachliche Humanvalidierung und der Nachweis der Produktionsreife. Die angemessene Gesamtbewertung lautet daher nicht, dass das System bereits verlässliche Versicherungsentscheidungen treffen kann, sondern dass es eine tragfähige technische und methodische Grundlage für ein kontrolliertes, quellengebundenes Assistenzsystem bereitstellt.

\section{Ausblick und zukünftige Arbeiten}

Die nächsten Schritte sollten sich an den beobachteten Fehlermechanismen und nicht ausschließlich an höheren Einzelmetriken orientieren. Kurzfristig sind die Überführung vorhandener Evidenz in vollständige Antworten, ein konsistentes Quellenhandling und die Absicherung der Zugriffsschicht zu priorisieren. Mittelfristig sind unabhängige Expertenbewertungen, wiederholte End-to-End-Läufe und produktionsnahe Lasttests erforderlich. Langfristig muss untersucht werden, ob sich die Architektur auf weitere Sprachen, Versicherungsbereiche und Organisationen übertragen lässt.

\subsection{Technische Weiterentwicklung}

Die sieben Fehlerfälle des aktuellen Referenzlaufs sollten als dauerhafte Regressionstests übernommen werden. Besonders wichtig ist eine eigenständige Vollständigkeitsprüfung neben der Groundedness-Prüfung. Dafür kann vor der Generierung aus Routingplan, CRM-Schema und Dokumentfrage eine strukturierte Liste erwarteter Aussagen und Evidenztypen gebildet werden. Nach der Generierung prüft eine zweite Stufe, ob jede Muss-Anforderung beantwortet, der richtigen Entität zugeordnet und bei dokumentbasierten Aussagen mit einer geeigneten Quelle verbunden wurde. Eine Antwort würde nur dann vollständig freigegeben, wenn Sicherheits-, Groundedness- und Vollständigkeits-Gate jeweils bestanden sind.

Als weiterführendes Architekturkonzept bietet sich ein maschinenlesbarer \emph{Evidenzvertrag} zwischen Routing, Datenzugriff, Generierung und Validierung an. Dieser Vertrag könnte für jede Anfrage zulässige Evidenzarten, erwartete Entitäten, Muss-Anforderungen, erlaubte Ableitungen und erforderliche Zitationen festhalten. Er würde die im Prototyp bereits vorhandene Herkunftstrennung zu einer expliziten, prüfbaren Schnittstelle erweitern. Dieses Konzept ist nicht Bestandteil der aktuellen Evaluation, ergibt sich aber unmittelbar aus dem beobachteten Verlust zwischen erfolgreichem Retrieval und unvollständiger Ausgabe.

Retrieval und Quellenverarbeitung sollten durch stabilere Dokument-, Seiten- und Abschnittsreferenzen verbessert werden, da Chunk-IDs nach einer Neuindexierung wechseln können. Zusätzlich sind Ablationen zu Chunk-Größe, Überlappung, Nachbarerweiterung, Tabellenverarbeitung, Kopf- und Fußzeilenbereinigung sowie Produktfiltern erforderlich. Eine OCR-Pipeline und die Verarbeitung gescannter Dokumente sollten erst nach eigener Evaluation als Systemfähigkeit ausgewiesen werden. Für die Laufzeit ist zunächst das Reranking zu optimieren, da es im aktuellen Lauf die langsamste instrumentierte Stufe war. Kandidatenzahl, Batch-Verarbeitung, Caching und alternative Modelle sollten dabei getrennt untersucht werden.

Vor einem realen Pilotbetrieb müssen Authentifizierung, rollenbasierte Autorisierung, TLS, Secret-Management und eine risikobasierte Fail-closed-Strategie umgesetzt werden. Für uneindeutige oder besonders folgenreiche Anfragen sollte das System nicht autonom mehr Werkzeuge nutzen, sondern gezielt Rückfragen stellen oder an eine Fachperson eskalieren. Ein solcher risiko- und evidenzadaptiver Kontrollpfad würde die bestehende deterministische Orchestrierung erweitern, ohne die Begrenzung der Werkzeugrechte aufzugeben.

\subsection{Erweiterung der Datengrundlage}

Das Dokumentkorpus sollte um weitere Produkte, Dokumentgenerationen, Versicherer und realistische Qualitätsprobleme erweitert werden. Dazu gehören gescannte Seiten, Tabellen, widersprüchliche Versionen, veraltete Bedingungen, fehlende Metadaten und ähnlich benannte Produkte. Ein zeitlicher Gültigkeitsbezug sollte Bestandteil der Referenzdaten werden, damit nicht nur semantische Relevanz, sondern auch die Auswahl der zum Anfragezeitpunkt gültigen Dokumentversion bewertet werden kann.

Für CRM-Tests sind größere synthetische Bestände mit fehlenden Werten, Dubletten, Namensgleichheiten, historischen Änderungen, fehlerhaften Beziehungen und unterschiedlichen Berechtigungsprofilen erforderlich. Kontakte, Policen und Schäden sollten jeweils eigene Selection-Labels erhalten. Damit ließe sich die Entitätsauswahl getrennt vom CRM Fact Recall messen und feststellen, ob eine fachlich korrekte Information aus dem richtigen Datensatz stammt.

Der balancierte 80-Fall-Datensatz sollte durch unabhängig formulierte Anfragen ergänzt werden, die reale Häufigkeiten, Tippfehler, Mehrfachanliegen, indirekte Referenzen und organisationsspezifische Abkürzungen abbilden. Reale oder realitätsnahe Daten dürfen erst nach einem dokumentierten Datenschutz- und Governance-Verfahren einbezogen werden. Für erste Erweiterungen sind synthetische oder konsequent anonymisierte Daten vorzuziehen.

\subsection{Weiterführende Evaluation}

Die höchste Priorität besitzt eine unabhängige Expertenannotation der aktuellen und zukünftigen End-to-End-Antworten. Bewertet werden sollten fachliche Korrektheit, Vollständigkeit, Entitäts- und Zahlenkonsistenz, Claim-Support, Citation Precision und Citation Recall. Mindestens zwei voneinander unabhängige Bewertende sowie ein dokumentiertes Verfahren zur Auflösung von Uneinigkeiten würden die Reliabilität der Labels erhöhen. Auf dieser Grundlage könnten Groundedness Precision, Recall und False Acceptance Rate für den tatsächlichen Anwendungskontext bestimmt werden.

Mehrere vollständige Wiederholungsläufe mit festgehaltenen Modell-, Prompt-, Korpus- und Konfigurationsständen sind erforderlich, um Run-to-Run-Variabilität zu quantifizieren. Kontrollierte Ablationen sollten anschließend den kausalen Beitrag von hybridem Retrieval, Reranking, deterministischer Policenfilterung, Groundedness, Vollständigkeitsprüfung und Guardrails untersuchen. Zusätzlich sind route-spezifische Last-, Parallelitäts-, Kaltstart-, Ausfall-, Kosten- und Langzeittests notwendig. Ergebnisse sollten weiterhin pro Systemstand und Verarbeitungspfad berichtet werden.

Eine Nutzerstudie mit Versicherungsfachkräften sollte schließlich untersuchen, ob das Assistenzsystem die Bearbeitungszeit, Fehlerquote und Nachvollziehbarkeit realer Informationsaufgaben verbessert. Neben objektiven Leistungsmaßen sind Vertrauen, Fehlinterpretationen, wahrgenommene Arbeitsbelastung und der Umgang mit Fallbacks zu erfassen. Erst die Verbindung technischer Metriken mit fachlichen Expertenurteilen und tatsächlichem Arbeitsnutzen erlaubt eine belastbare Aussage über praktische Wirksamkeit.

\subsection{Übertragung auf weitere Sprachen und Versicherungsbereiche}

Die sprachliche Übertragbarkeit sollte zunächst mit einem vollständig deutschsprachigen Dokumentkorpus, deutschen Referenzantworten und natürlich formulierten deutschen Anfragen geprüft werden. Das mehrsprachige Embedding-Modell allein belegt keine gleichwertige Retrieval- und Antwortqualität. Routingregeln, Fachbegriffe, Komposita, Abkürzungen und Zitationsanforderungen müssen sprachspezifisch getestet und gegebenenfalls neu kalibriert werden.

Danach kann die Architektur auf weitere Versicherungssparten, Dokumenttypen und CRM-Schemata übertragen werden. Jede Erweiterung erfordert neue Produktregeln, Evidenztypen, Referenzfälle und Berechtigungsmodelle. Eine Übertragung auf andere regulierte Domänen ist konzeptionell denkbar, setzt aber eine domänenspezifische Definition zulässiger Werkzeuge, Freigabekriterien, Eskalationswege und menschlicher Verantwortlichkeiten voraus.

\subsection{Weiterführende Forschungsfragen}

Aus der Arbeit ergeben sich insbesondere folgende weiterführende Forschungsfragen:

\begin{enumerate}
    \item Wie kann die Vollständigkeit versicherungsfachlicher Antworten unabhängig von ihrer Groundedness zuverlässig und reproduzierbar geprüft werden?
    \item In welchem Maß reduziert ein maschinenlesbarer Evidenzvertrag Fehler bei Entitätsbindung, Muss-Anforderungen und Zitationen, ohne die Zahl unnötiger Fallbacks zu erhöhen?
    \item Welche Kombination aus deterministischem Routing, Unsicherheitsschätzung und gezielten Rückfragen behandelt mehrdeutige oder mehrteilige Anfragen am sichersten?
    \item Wie stabil bleiben Pfadauswahl, Antwortqualität und Quellenzuordnung über wiederholte Läufe, Modellversionen, Sprachen und veränderte Dokumentbestände hinweg?
    \item Unter welchen organisatorischen und technischen Bedingungen verbessert ein quellengebundenes CRM--RAG-System nachweisbar die Arbeitsqualität und Bearbeitungszeit von Versicherungsfachkräften?
\end{enumerate}

Die vorliegende Arbeit schafft eine Grundlage für diese Untersuchungen. Ihr zentraler Befund ist, dass eine kontrollierte agentische CRM--RAG-Architektur strukturierte und unstrukturierte Versicherungsinformationen technisch zuverlässig zusammenführen kann, wenn Pfade, Werkzeuge, Evidenzherkunft und Freigabekriterien explizit begrenzt werden. Zugleich zeigt die Evaluation, dass technische Ausführbarkeit nicht mit fachlicher Vollständigkeit oder Produktionsreife gleichzusetzen ist. Der nächste Fortschritt liegt daher weniger in zusätzlicher Autonomie als in präziseren Evidenzverträgen, unabhängiger Validierung und einer konsequenten Trennung von Betriebs-, Qualitäts- und Sicherheitsfreigabe.
